\documentclass[a4paper]{article}
\usepackage[polish]{babel}
\usepackage[cp1250]{inputenc}
%\usepackage{polski}
\usepackage{amsfonts}
\usepackage[T1]{fontenc}
\usepackage[dvips]{graphicx}
%\usepackage{pldoc}

\pagestyle{headings}
\textwidth      15.5cm
\oddsidemargin    .1cm
\evensidemargin   .1cm

\begin{document}

\thispagestyle{empty}

\begin{minipage}{5cm}
\includegraphics[scale=0.7]{salomon.eps}
\end{minipage}
\begin{minipage}{10cm}
\begin{center}
{\Large Salomon\\
System przetwarzania wiedzy\\}
\end{center}
\end{minipage}

\vspace*{0.5cm}

\hrulefill

\vspace*{1cm}


\begin{flushright}
{\Large \emph{Vision Scope, ver 1.0} }
\end{flushright}

\vspace{7cm}

\begin{flushleft}
{ \Large Historia wersji }\[\]
\begin{tabular}{|p{3.5cm}||p{3.5cm}|p{3.5cm}|p{3.5cm}|}
  \hline
  Data & Wersja & Opis & Autor \\
  \hline \hline
  17-04-2005 & 1.0 & Initial version & Jakub Ga�kowski \\
  \hline
\end{tabular}
\end{flushleft}

%\[\] \[\] \[\] \[\] \[\] \[\] \[\] \[\] \[\]

%\begin{flushleft}
%{ \Large Autor} \[\]
%\begin{tabular}{|p{3.5cm}|p{3.5cm}|p{4.5cm}|p{3.5cm}|}
%  \hline 
%  Jakub Ga�kowski & Developer & avi@student.uci.agh.edu.pl & 17-04-2005 \\
%  \hline
%\end{tabular}
%\end{flushleft}

\newpage

\section{Opis produktu}
Ostateczna wersj� produktu b�dzie system wyposa�ony w pe�n� obs�ug� zbior�w danych, regu� i atrybut�w oraz mo�liwo�� pracy z danymi zawartymi w r�nych bazach danych wraz z kompletn� dokumentacj� techniczn� i u�ytkow�.

\section{Wersje programu}
\begin{itemize}

\item Wersja 1.0
	\begin{itemize}
		\item stabilna wersja Salomona
    \item poprawna obs�uga zbior�w danych
    \item prezentacja dzia�ania prostych algorytm�w opartych o zbiory danych (tworzenie zbior�w danych w oparciu o zadane kryteria, podstawowe operacje na zbiorach w odniesieniu do zbior�w danych - suma, r�nica, przeci�cie)
	\end{itemize}

\item Wersja 1.1
	\begin{itemize}
    \item zaawansowana obs�uga zbior�w danych (generowanie nowych zbior�w danych na podstawie istniej�cych)
    \item prezentacja rozbudowanych algorytm�w operuj�cych na zbiorach danych    
	\end{itemize}

\item Wersja 1.2 
	\begin{itemize}
    \item podstawowa obs�uga wiedzy opartej o atrybuty i regu�y
    \item prezentacja dzia�ania prostych algorytm�w opartych o regu�y i atrybuty (tworzenie regu� na podstawie zadanych kryteri�w)
	\end{itemize}

\end{itemize}

\end{document}


